\documentclass[11pt, ngerman]{scrartcl}
\usepackage[utf8]{inputenc}

\usepackage{scrlayer-scrpage}
\pagestyle{scrheadings}\usepackage{scrlayer-scrpage}
\pagestyle{scrheadings}

\renewcommand{\headfont}{\normalfont}

\ihead{\small Seminarvortrag}
 \chead{  \hspace{1,3cm }\small Universalität des homogenen Wiener Chaos}
 \ohead{\small SoSe 2025}
 \setheadsepline{.2pt}


% \usepackage[notcite,notref]{showkeys}
\setlength{\parindent}{0cm}
\usepackage{amsthm}
\usepackage{amsmath, amssymb, amsfonts}
\renewcommand{\proofname}{Beweis}
\theoremstyle{definition}
\newtheorem{definition}{Definition}
\newtheorem{beispiel}[definition]{Beispiel}
\newtheorem{bemerkung}[definition]{Bemerkung}
\newtheorem{Fdefinition}[subsection]{Definition}
\newtheorem{Aufgabe}{Übung}
\theoremstyle{plain}
\newtheorem{lemma}[definition]{Lemma}
\newtheorem{satz}[definition]{Satz}
\newtheorem{hauptsatz}[definition]{Hauptsatz}
\newtheorem{corollar}[definition]{Korollar}
\newtheorem{proposition}[definition]{Proposition}
\newtheorem{folgerung}[definition]{Folgerung}
\newtheorem{theorem}[definition]{Theorem}
\newtheorem{Axiom}[definition]{Axiom}
\newtheorem{Abschaetzung}[definition]{Abschätzung}

% Jetzt zu den Abkürzungen
% alle Mathe-Abkürzungen
\renewcommand{\epsilon}{\varepsilon} % wichtige Sache
\newcommand{\loc}[3]{L_{#1}^{#2}(#3)}
\newcommand{\s}[0]{\mathfrak{s}}
\newcommand{\m}[0]{\mathfrak{m}}
\newcommand{\br}[1]{\left(#1\right)}
\newcommand{\lbr}[1]{\left(#1\right}
\newcommand{\C}[0]{\mathbb{C}}
\newcommand{\R}[0]{\mathbb{R}}
\newcommand{\Q}[0]{\mathbb{Q}}
\newcommand{\M}[0]{\mathbb{M}}
\newcommand{\Z}[0]{\mathbb{Z}}
\newcommand{\N}[0]{\mathbb{N}}
\newcommand{\X}[0]{\mathbb{X}}
\newcommand{\F}[0]{\mathcal{F}}
\newcommand{\Fi}[0]{\mathbb{F}}
\newcommand{\G}[0]{\mathbb{G}}
\newcommand{\W}[0]{\mathbb{W}}
\newcommand{\B}[0]{\mathcal{B}}
\renewcommand{\H}[0]{\mathcal{H}}
\newcommand{\dP}[0]{\mathrm{d}P}
\newcommand{\dPdif}[1]{\mathrm{d}#1}
\renewcommand{\S}{\textcolor{orange}{\textbf{(S)}}\,}
\newcommand{\homQG}[2]{Q_{d_{#1}}(N_{#2}^{(#1)}, f_{#2}^{(#1)}, \mathbb{G})}
\newcommand{\homQGo}[1]{Q_{d_{#1}}(N^{(#1)}, f^{(#1)}, \mathbb{G})}
\newcommand{\homQXo}[1]{Q_{d_{#1}}(N^{(#1)}, f^{(#1)}, \mathbb{X})}
\newcommand{\homQX}[2]{Q_{d_{#1}}(N_{#2}^{(#1)}, f_{#2}^{(#1)}, \mathbb{X})}
\newcommand{\Ch}[1]{\mathcal{H}_{#1}}
\newcommand{\Pol}[1]{\mathcal{P}_{#1}}
\newcommand{\Hil}[0]{\mathfrak{h}}
\usepackage{bbm} %for the indicator
\newcommand{\indicator}[1]{\mathbbm{1}_{#1}}
\newcommand{\indicatorset}[1]{\mathbbm{1}_{\{#1\}}}
\newcommand{\law}[0]{\overset{\mathcal{L}}{=}}
\newcommand{\normald}[2]{\mathcal{N}\left(#1, #2\right)}
\newcommand{\E}[1][]{%
    \operatorname{E}%
    \ifx\\#1\\\else% Check if optional argument is provided
        \left[#1\right]%
    \fi%
}
\renewcommand{\P}[1][]{%
    \operatorname{P}%
    \ifx\\#1\\\else% Check if optional argument is provided
        \left(#1\right)%
    \fi%
}
\newcommand{\<}{\langle}
\renewcommand{\>}{\rangle}
\newcommand{\Cov}[0]{\operatorname{Cov}}
\newcommand{\iid}[0]{i.\,i.\,d.\,}
\renewcommand{\P}[0]{\operatorname{P}}
\newcommand{\Czi}[0]{\mathcal{C}([0, \infty))}
\newcommand{\brp}[1]{\langle{}#1\rangle{}}

\newcommand{\norm}[1]{\lVert#1\rVert}
\usepackage{aligned-overset}
\newcommand\restr[2]{\ensuremath{\left.#1\right|_{#2}}} % für restr-Abbildung: \restr{f}{A}


\begin{document}
\begin{center}

\vspace{0,25cm}

    \large{\textbf{Universalität des homogenen Wiener Chaos}}\\
\end{center}
\vspace{0,25cm}
\textbf{Struktur des Vortrags}
\begin{itemize}
    \item Wir stellen das Setting vor.
    \item Wir präsentieren das Hauptresultat und was es mit dem Wiener Chaos zu tun hat.
    \item Wir führen \textbf{Kontraktionen und Einflussindizes} ein -- die \textit{Stars} des Beweises.
    \item Mithilfe einiger erster Techniken legen wir die Beweisidee dar. 
    \item Wir können für die Konvergenz mittels Dreiecksungleichung zwei Terme gesondert betrachten. Dies tuen wir dann. 
    \item Am Ende geben wir skizzenhaft wieder, wie man den Beweis vervollständigt.
    \end{itemize}

\begin{definition}[Homogene Summen]
 Sei $\mathbb{X} = (X_n)_{n \in \N}$ eine Folge von unabhängigen zentrierten Zufallsvariablen mit Varianz 1 (mit gewissen Momentannahmen). 
        Für $d \in \N$ und $f^{(N)}: \{1, \ldots, N\}^d \to \R$ symmetrisch mit $f^{(N)}(i_1, \ldots, i_d) = 0$ falls $i_j = i_k$ für $j \neq k$ definieren wir die homogene Summe
        \[
        Q_d(f^{(N)}, N, \mathbb{X}) = \sum_{1 \leq i_1, \ldots, i_d \leq N} f^{(N)}(i_1, \ldots, i_d) X_{i_1} \cdots X_{i_d}.
        \]
\end{definition}
\begin{definition}[$\mathcal{H}$-Abstand und Kolmogorov-Abstand]
        Sei $\mathcal{H}$ eine Klasse von Funktionen $h: \R^m \to \R$, so definiere
            $d_{\mathcal{H}}(F, Z) = \sup_{h \in \mathcal{H}} \left|\E[h(F)]- \E[h(Z)]\right|.$

        Der Kolmogorov-Abstand zwischen zwei multivariaten Verteilungen bzw. \textbf{Verteilungsfunktionen} $F$ und $G$ ist gegeben durch $
        d_{\text{Kol}}(F,G) = \sup_{x \in \R^m} |F(x_1, \ldots, x_m) - G(x_1, \ldots, x_m)|,$
        also 
        \[
            d_{\text{Kol}}(F,G) = d_{\mathcal{H}}(F,G), \mathcal{H} = \{\indicator{(-\infty, x_1]\times \cdots \times (-\infty, x_m]}\}_{x_1, \ldots, x_m \in \R}.
        \]
    \end{definition}
\begin{theorem}[Universalität des Wiener Chaos]
        Sei $\mathbb{G} = (G_n)_{n \in \N}$ eine Folge von unabhängigen standardnormalverteilten Zufallsvariablen. 
        Fixiere $m \geq 1$ und $d_1, \ldots, d_m \geq 2$. Sei ferner für jedes $j = 1, \ldots, m \; \{(N_n^{(j)}, f_n^{j})\}_{n \geq 1}$ so, dass $N_n^{(j)} \overset{n \to \infty}{\longrightarrow} \infty$ sei und $f_n^{(j)}: [N_n^{(j)}]^{d_j} \to \R$ auf Diagonalen ($i_k = i_j , j \neq k$) verschwinde. 

        Seien nun für jedes $j = 1, \ldots, m$ die Erwartungswerte $\E[Q_{d_j}(N_n^{(j)}, f_n^{(j)}, \mathbb{G})^2]$ beschränkt. 
        Sei $\Sigma$ eine positiv semidefinite symmetrische $m\times m$-Matrix mit nicht-verschwindenden Diagonaleinträgen. Nenne dieses Setting \textbf{(S)}.

        Dann sind für $n \to \infty$ äquivalent: 
        \begin{itemize}
            \item[(1)] $\left(Q_{d_j}(N_n^{(j)}, f_n^{(j)}, \mathbb{G})\right)_{j=1}^m \Longrightarrow \mathcal{N}(0, \Sigma)$
            \item[(2)]  Für jede Folge $\mathbb{X} = \{X_i\}_{i \in \N}$ unabhängiger zentrierter Zufallsvariablen mit Varianz 1 so, dass $\sup_{i} \E[|X_i|^3] < \infty$, konvergiert \[
                \left(Q_{d_j}(N_n^{(j)}, f_n^{(j)}, \mathbb{X})\right)_{j=1}^m \overset{d_{\text{Kol}}}{\longrightarrow} \mathcal{N}(0, \Sigma).
            \]
        \end{itemize}
    \end{theorem}
\begin{definition}[Kontraktionen und Einflussindizes]
        Seien $d_1, d_2 \geq 2$ und $N\in \N \cup \{\infty\}$ fixiert und es sei $f: [N]^{d_1} \to \R, g: [N]^{d_2} \to \R$ symmetrische Funktionen, die auf Diagonalen verschwinden. 
        Definiere für jedes $r = 0, \ldots, d_1 \wedge d_2$ die Kontraktion $f *_r g: [N]^{d_1 + d_2 - 2r} \to \R$ durch \begin{align*}
        f *_r& g(i_1, \ldots,  i_{d_1 + d_2 - 2r}) = \\ &\sum_{1 \leq a_1, \ldots, a_r \leq N} f(a_{1}, \ldots, a_r, i_1, \ldots, i_{d_1 -r}) g(a_1, \ldots, a_r, i_{d_1 - r +1}, \ldots, i_{d_1 + d_2 - 2r}).
        \end{align*}
        Definiere ferner $\norm{f}_d^2 := \sum_{1 \leq i_1, \ldots, i_d \leq N} f^2(i_1, \ldots, i_d)$.
        Der $i$te Einflussindex von $f$ ist gegeben durch \[
        \text{Inf}_i(f) = \sum_{\{i_2, \ldots, i_d\} \subset [N]^{d-1}} f^2(i, i_2, \ldots, i_d) = (d-1)!^{-1} \sum_{1 \leq i_2, \ldots, i_d \leq N} f^2(i, i_2 \ldots, i_d).
        \]
    \end{definition}
\begin{definition}[Wiener-Chaos-Notation]
    Sei $(\Hil, \<\cdot, \cdot\>_{\Hil})$ ein separabler Hilbertraum und $\{e_i\}_{i \in \N}$ eine Orthonormalbasis von $\Hil$. Es bezeichnen $\Hil^{\otimes k}$ und $\Hil^{\odot k}$ das $k$-fache (resp. $k$-fache symmetrische) Tensorprodukt von $\Hil$. 
    Es bezeichne $\zeta$ einen isonormalen gaußschen Prozess auf $\Hil$. 
    Mit $I_k(h)$ bezeichne für $h \in \Hil^{\otimes k}$ $k$-faches Wiener-Integral bezüglich $\zeta$ und definiere \[\Ch{k} = \{I_k(h)\}_{h \in \Hil^{\otimes k}}, \Pol{k} = \overline{\<\{\text{Polynome von Grad } \leq k \text{ in } \zeta h\}\>}\]
    \end{definition}
    \begin{itemize}
        \item Verbindung zu unserem Setting: Für $e_1, \ldots, e_k$ orthonormal gilt \[I_k(e_1 \otimes \ldots \otimes e_k) = I_1(e_1) \cdots I_1(e_k). 
        \]
        \[
        h = \sum_{1 \leq i_1, \ldots, i_d \leq N} f(i_1, \ldots, i_d) e_{i_1} \otimes \cdots \otimes e_{i_d} \in \Hil^{\odot d}
        \] rechtfertigt somit $Q_d(N,f,\mathbb{G}) = I_d(h)$.
    \end{itemize}

\begin{Abschaetzung}
    \begin{align*}
        &\norm{f_n^{(j)} *_{d_j -1 } f_n^{(j)}}_{2}^2  \geq  
         \sum_{1 \leq i \leq N} \left[\sum_{1 \leq i_2, \ldots, i_d \leq N} (f_n^{(j)})^2(i, i_2, \ldots, i_d)\right]^2 \\ & \qquad \geq \max_{i = 1, \ldots, N} \left[\sum_{1 \leq i_2, \ldots, i_d \leq N} (f_n^{(j)})^2(i, i_2, \ldots, i_d)\right]^2 \geq \max_{i = 1, \ldots, N}  \left[(d-1)! \text{Inf}_i(f_n^{(j)}) \right]^2
        \end{align*}
\end{Abschaetzung}
\begin{satz}[Abstand von Chaos und verrauschter homogener Summe]
       Seien $\mathbb{X}, \mathbb{G}$ wie immer, fixiere $m \geq 1; d_m \geq \cdots \geq d_1 \geq 2; N_1, \ldots, N_m \geq 1$ und $f_j: [N_j]^{d_j} \to \R$ seien symmetrisch und 0 auf Diagonalen. 
       Setze $\E[\homQGo{j}^2] = 1$ (und somit $\E[\homQXo{j}^2] =1$) für alle $j = 1, \ldots, m$ voraus.
        Nimm ferner an, es gibt ein $C > 0$ derart, dass $\sum_{i=1}^{\max_{j}N_j} \max_{1 \leq j \leq m} \mathrm{Inf}_i(f_j) \leq C.$ Dann gilt für all dreifach differenzierbaren $\varphi: \R^m \to \R$ mit $\norm{\varphi'''}_\infty < \infty$ und $\varphi, \varphi', \varphi'' \in L^2(P^{\mathbf{Q}(\Tilde{\mathbb{X}})})$ für alle $\Tilde{\mathbb{X}}$, dass \begin{align*}
        &\left|\E[\varphi(Q^{(1)}(\mathbb{X}), \ldots, Q^{(m)}(\mathbb{X}))] - \E[\varphi(Q^{(1)}(\mathbb{G}), \ldots, Q^{(m)}(\mathbb{G}))]\right| \\ \leq &  C \cdot \norm{\varphi'''}_\infty \left(\beta + \sqrt{\frac{8}{\pi}}\right)  \times \left[\sum_{j=1}^m (16 \sqrt{2} \beta)^{\frac{d_j-1}{3}} d_j!\right]^3 \sqrt{\max_{j=1, \ldots, m}\max_{1 \leq i \leq N} \mathrm{Inf}_i(f_j)}.
        \end{align*}
        Dabei ist $\beta := \sup_{i \in \N} \E[|X_i|^3] <\infty$ und $\norm{\varphi'''}_{\infty} := \max_{|\alpha| = 3} \sup_{x \in \R^m} |\partial^\alpha \varphi(x)|$, $\norm{\varphi'}_2, \norm{\varphi''}_2$ analog
         ($\alpha \in \N^{m}$) \end{satz}

\begin{definition}[Wesentliche Bestandteile des Beweises]
  \[
        \mathbb{Z}^{(i)} = \left(G_1, \ldots, G_i, X_{i+1}, \ldots, X_{\max_j N_j}\right).\]
  \begin{align*}
        U_j^{(i)} &:= \sum_{\substack{1 \leq i_1, \ldots, i_{d_j} \leq N_j \\ \forall k: i_k \neq i}} f_j(i_1, \ldots, i_{d_j}) Z_{i_1}^{(i)} \cdots Z_{d_j}^{(i)} \text{ für } j = 1, \ldots, m\\
        V_j^{(i)} &:= \sum_{\substack{1 \leq i_1, \ldots, i_{d_j} \leq N_j \\ \exists k: i_k = i}} f_j(i_1, \ldots, i_{d_j}) Z_{i_1}^{(i)} \cdots \widehat{Z_{i_{k}}^{(i)}} \cdots Z_{i_d}^{(i)} \text{ für } j = 1, \ldots, m
    \end{align*}
\end{definition}
\begin{itemize}
  \item Zentraler Punkt: Wegen Zentriertheit, Unabhängigkeit und Normierheit \begin{align*}\E[(V_j^{(i)})^2] &= \sum_{\substack{1 \leq i_1, \ldots, i_{d_j} \leq N_j \\ \exists k: i_k = i}} f_j(i_1, \ldots, i_{d_j})^2 \E[(Z_{i_1}^{(i)})^2] \cdots \widehat{1}_{(i)} \cdots \E[(Z_{i_d}^{(i)})^2] \\ &= \sum_{\substack{1 \leq i_1, \ldots, i_{d_j} \leq N_j \\ \exists k: i_k = i}} f_j(i_1, \ldots, i_{d_j})^2 = d_j! \mathrm{Inf}_i(f_j) \end{align*} 
\end{itemize}

 \begin{satz}{Abstand von verrauschter Summe und der Normalverteilung}
        Wir befinden uns im Setting des vorherigen Satzes und nehmen zusätzlich an, dass $Z \sim \mathcal{N}(0, \Sigma)$, wobei $\Sigma$ wie zuvor -- und gleichzeitig noch die Kovarianzmatrix von $\mathbf{Q}(\mathbb{G})$ -- sei.
        Dann gilt für alle $\varphi: \R^m \to \R$ mit $\norm{\varphi'''}_\infty, \norm{\varphi''}_\infty < \infty$ und $\varphi, \varphi', \in L^2(P^{\mathbf{Q}(\Tilde{\mathbb{X}})})$ für alle $\Tilde{\mathbb{X}}$, dass  (für $D = D(\beta, \max_j d_j, m)$ von vorher)
        \vspace{-0.4cm}
        \begin{align*}
            &\left|\E[\varphi(\mathbf{Q}(\X))] - \E[\varphi(Z)]\right| \leq \norm{\varphi'''}_\infty \cdot D \sqrt{\max_{j=1, \ldots, m}\max_{1 \leq i \leq N} \mathrm{Inf}_i(f_j)} \\ &\qquad \qquad + \norm{\varphi''}_\infty \cdot \left(\sum_{i=1}^m \triangle_{ii} + 2 \cdot \sum_{1\leq i < j \leq m} \triangle_{ij}\right)\text{, wobei } \\
       & \triangle_{i,j} := \frac{d_j}{\sqrt{2}} \sum_{r = 1}^{d_j -1} (r-1)! \binom{d_i-1}{r-1} \binom{d_j-1}{r-1} \\ & \qquad \qquad \times  \sqrt{(d_i + d_j - 2r)!} \left(\norm{f_i *_{d_i -r } f_i}_{2r} + \norm{f_j *_{d_j - r}f_j}_{2r}\right) \\ & \qquad \qquad \qquad \qquad \qquad + \indicatorset{d_i < d_j} \sqrt{d_j! \binom{d_j}{d_i} \norm{f_j *_{d_j - d_i} f_j}_{2d_i}} \end{align*}
    \end{satz}

    \begin{definition}[Kontraktionen in symmetrischen Tensorprodukten]
            Sei $\Hil$ ein separabler Hilbertraum und $\{e_i\}_{i \in \N}$ eine Orthonormalbasis von $\Hil$. Dann definiere für $h = \sum_{i_1, \ldots, i_d} a(i_1, \ldots, i_d) e_{i_1} \otimes \cdots \otimes e_{i_d} \in \Hil^{\odot d_1}$ und und $g = \sum_{i_1, \ldots, i_d} b(i_1, \ldots, i_d) e_{i_1} \otimes \cdots \otimes e_{i_d} \in \Hil^{\odot d_2}$ für $p \in \{1, \ldots, d_1 \wedge d_2\}$ die $p$te Kontraktion $h \otimes_p g \in \Hil^{d_1 + d_2 - 2p}$  als \[
            h \otimes_p g := \sum_{\substack{i_1, \ldots, i_{d_1 - p} \\ j_1, \ldots, j_{d_2 - p}}} a *_p b(i_1, \ldots, j_{d_2 - p}) e_{i_1} \otimes \cdots \otimes e_{i_{d_1 - p}} \otimes e_{j_1} \otimes \cdots \otimes e_{j_{d_2 - p}}.
            \]
            \begin{itemize}
            % \item Beachte: Man kann schreiben \begin{align*}
            %  &h \otimes_p g = \<h,g\>_{\Hil^{\odot p}}  \\ &= \sum_{\substack{i_1, \ldots, i_{d_1 - p} \\ j_1, \ldots, j_{d_2 - p}}} \big\<\sum_{i_{d_1-p+1, \ldots, i_{d_1}}} a(i_{1}, \ldots, i_{d_1}) e_{i_{d_1-p+1}} \otimes \cdots \otimes e_{i_{d_1}} ,  \\ &  \sum_{i_{j_{d_2-p+1}, \ldots, j_{d_2}}} b(j_{1}, \ldots, j_{d_2}) e_{j_{d_2-p+1}} \otimes \cdots \otimes e_{i_{d_2}}\big\>_{\Hil^{\odot p}} e_{i_1} \otimes \cdots \otimes e_{i_{d_1 - p}} \otimes e_{j_1} \otimes \cdots \otimes e_{j_{d_2 - p}}.
            % \end{align*}
            \item Insbesondere für $\{e_i\} \subset \ell^2(\N)$ gegeben durch $e_i(n) = \delta_{i,n}$ (Kronecker-Delta) haben wir gerade \[f *_r g = f \otimes_r g \in \ell^2(\N)^{\otimes d_1 + d_2 - 2r}.\]
            \end{itemize}
        \end{definition}
        
  \begin{satz}[Multiplikationsformel für Wienersche Mehrfachintegrale]
    Seien $f \in \Hil^{\odot p}, g \in \Hil^{\odot q}$. Dann gilt $I_p(f)I_q(g) = \sum_{r=0}^{p \wedge q} r! \binom{p}{r}\binom{q}{r}I_{p+q-2r}(f \otimes_r g)$.
  \end{satz}
\end{document}